\chapter{求解算法设计}

\section{总体流程}
算法按时间推进，每次处理一个滑动窗口，流程如下：
\begin{enumerate}
  \item 读取当前窗口 RSSI 数据并做异常预处理；
  \item 基于第 4 章方法构造初始轨迹；
  \item 执行迭代重加权最小二乘求解；
  \item 输出当前时刻位置并滑动到下一窗口。
\end{enumerate}

\section{IRWLS 主算法}
设窗口变量为 $\bm X$，在第 $k$ 次迭代：
\begin{equation}
\bm X^{(k+1)}=\bm X^{(k)}+\Delta\bm X^{(k)},
\end{equation}
其中 $\Delta\bm X^{(k)}$ 由线性化法方程求解得到。

\begin{algorithm}[H]
\caption{滑动窗口鲁棒定位算法}
\begin{algorithmic}[1]
\State 输入：窗口观测、初始轨迹、参数 $\lambda_s,\lambda_h,\delta$
\For{$k=1$ to MaxIter}
\State 计算各观测残差 $r_i$
\State 根据 Huber 规则更新权重 $\omega_i$
\State 构造线性系统 $\bm A\Delta\bm X=\bm b$
\State 求解增量并更新 $\bm X\leftarrow\bm X+\Delta\bm X$
\State 执行可行域投影修正
\If{$\|\Delta\bm X\|/\|\bm X\|<\varepsilon$}
\State break
\EndIf
\EndFor
\State 输出窗口末端位置估计
\end{algorithmic}
\end{algorithm}

\section{参数设置原则}
核心参数包括：
\begin{itemize}
  \item $\lambda_s$：轨迹平滑强度，过大将抑制转弯；
  \item $\lambda_h$：方向变化惩罚，控制航向突变；
  \item $\delta$：Huber 阈值，决定异常点降权力度；
  \item $w$：窗口长度，影响时延与稳定性权衡。
\end{itemize}

\section{收敛与复杂度}
窗口维度为 $2w$，每步求解复杂度约为 $\mathcal O(w^3)$（稠密求解）或更低（稀疏求解）。在课程实验配置 $w=8$ 下，单帧平均计算时间可控制在毫秒级。

\section{失败保护机制}
为避免极端场景下发散，设置以下保护：
\begin{enumerate}
  \item 迭代步数上限；
  \item 增量过大时执行阻尼缩放；
  \item 连续异常时回退到预测轨迹；
  \item 信标数量不足时启用短时外推模式。
\end{enumerate}
